%% ===================================================================
\section{Introduction}\label{sec:intro}
%% ===================================================================

Inclusive production of a single gluon in the collision of a dilute projectile with a dense target is the simplest observable that probes the nonlinear regime of QCD at high energy. In the colour glass condensate (CGC) effective theory~\cite{McLerran:1993ni,McLerran:1993ka} the target is described by eikonal Wilson lines $U(x)$ in the adjoint representation, and the projectile by the colour charge density $\rho^a(x)$ of its fast (valence) partons together with the cloud of soft gluons that these charges radiate. At leading order (LO) in the strong coupling the gluon spectrum is known in closed form~\cite{Kovchegov:1998bi,Kovchegov:2001sc,Dumitru:2001ux,Blaizot:2004wu,Kovner:2006wr}. The measured gluon is emitted by a valence charge either before or after the eikonal interaction with the target. The cross section is a convolution of the projectile charge correlator $\langle\rho\rho\rangle$ with a target correlator of the form $\Tr\{[\Ud(x)-\Ud(w)][U(x')-U(w')]\}$.

At high energy these correlators are dressed by large logarithms of the energy. At leading-logarithmic (LL) accuracy the produced gluon at rapidity $y$ divides the evolution into two parts~\cite{Kovchegov:2001sc,Baier:2005dv,Kovner:2006ge,Kovner:2006wr,Gelis:2008rw}. The projectile is evolved from its own rapidity down to $y$ and the target from its rapidity up to $y$. In the dilute projectile the charge correlator then obeys the BFKL equation~\cite{Kuraev:1977fs,Balitsky:1978ic}. The target Wilson lines obey the JIMWLK equation~\cite{JalilianMarian:1997jx,JalilianMarian:1997gr,JalilianMarian:1998cb,Kovner:2000pt,Weigert:2000gi,Iancu:2000hn,Ferreiro:2001qy}, whose large-$N_c$ mean-field limit is the Balitsky--Kovchegov (BK) equation~\cite{Balitsky:1995ub,Kovchegov:1999yj}.

At next-to-leading order (NLO) much less is known about this observable. The evolution equations themselves are available at NLO: the NLO BK equation~\cite{Balitsky:2007feb} and the NLO JIMWLK Hamiltonian~\cite{Balitsky:2013fea,Kovner:2013ona,Kovner:2014lca,Lublinsky:2016meo}, together with their running-coupling content~\cite{Balitsky:2006wa,Kovchegov:2006vj}. Running-coupling corrections to the LO gluon spectrum were obtained in Ref.~\cite{Horowitz:2010yg}. Recently it was shown that some of the $\beta_0$-dependent transverse logarithms of NLO JIMWLK are not related to the running of the coupling. They are instead the first step of DGLAP evolution of the projectile~\cite{Kovner:2023vsy,Altinoluk:2023hfz}. Complete NLO results exist for single inclusive \emph{hadron} production at forward rapidity in the hybrid formalism~\cite{Chirilli:2011km,Chirilli:2012jd,Stasto:2013cha,Altinoluk:2014eka,Ducloue:2016shw,Iancu:2016vyg}. There the projectile is described by collinear parton distributions, and only the target undergoes small-$x$ evolution. These studies showed that the subtraction of the rapidity divergence, and the precise rapidity interval assigned to the evolution, are delicate at NLO.

At mid rapidity both the projectile and the target are probed at small $x$, and both must be treated in the same framework. In this paper we compute the single inclusive gluon spectrum at NLO, i.e.\ to order $g^4$, in the light-cone wave function (LCWF) approach~\cite{Lepage:1980fj,Brodsky:1997de,Kovner:2005nq,Kovner:2013ona,Kovner:2014lca,Lublinsky:2016meo}. The soft gluon modes of the projectile, with longitudinal momenta $\Lambda<k^+<\vee$, are dressed by a unitary operator $\Omega$ that acts on the valence state. The coefficients of $\Omega$ are fixed to order $g^3$ by light-cone perturbation theory (LCPT). The cross section is then the expectation value of the number operator of \emph{dressed} gluons in the scattered state. This construction treats real and virtual corrections, and initial- and final-state radiation, on the same footing.

Our main results are the following.
\begin{enumerate}
\item We give the NLO LCWF of the valence vacuum to order $g^3$ and of a single-gluon state to order $g^2$ in transverse coordinate space. We use them to fix the coefficients $\mf N$, $\mf A^{(1)}$, $\mf A^{(3)}$, $\mf B_2$, $\mf B_3$ and $\mf C_1$ of $\Omega$; the remaining coefficients follow from unitarity (Secs.~\ref{sec:lcwf} and~\ref{sec:omega}). The one-gluon wave function has two orderings, with the outgoing soft gluon harder or softer than the incoming one. The relative sign of their non-instantaneous parts is fixed by requiring that the double pole $1/(p^+-k^+)^2$ of the instantaneous interaction cancels on both sides of $p^+=k^+$.
\item We write the NLO cross section as a sum of products of these coefficients. There are five classes of virtual (one-gluon) terms, four classes of squared real (two-gluon) amplitudes and six classes of interference terms (Sec.~\ref{sec:xsec}).
\item At LL accuracy the full rapidity interval $\delta Y=\log(\vee/\Lambda)$ opened by the unobserved gluon is shared between the projectile and the target (Sec.~\ref{sec:LL}):
\begin{equation}
\frac{d^3N}{d^3k}\bigg|_{\rm LL}=\log\frac{\vee}{k^+}\,\big[{\rm BFKL}\otimes{\rm LO}\big]+\log\frac{k^+}{\Lambda}\,\big[{\rm JIMWLK}\otimes{\rm LO}\big].
\label{eq:intro-main}
\end{equation}
Real and virtual terms combine into a double commutator of soft-gluon emission operators. Its action on the LO cross section is exactly that of the JIMWLK Hamiltonian, including the terms in which the soft gluon is emitted by the measured gluon itself. All contributions with three or four valence charges cancel.
\item We define the NLO remainder (the ``finite part'') by an exact subtraction of $\log(k^+/\Lambda)$ and $\log(\vee/k^+)$, keeping all dependence on the finite upper cutoff $\vee$. We give this remainder term by term and show that all power-like terms $\propto1/\Lambda$ cancel (Sec.~\ref{sec:finite}).
\item The real gluon splitting function appears in four places: in the self-energy of the emitted gluon, in the loop across the target, and in the initial-state, interference and final-state collinear splittings. The real collinear terms come with relative weights $1:-2:1$. We use this to separate the $\beta_0$ terms associated with the running coupling from those associated with DGLAP evolution of the projectile and with fragmentation, along the lines of Ref.~\cite{Kovner:2023vsy} (Sec.~\ref{sec:rc}).
\end{enumerate}

The paper is organized as follows. Section~\ref{sec:framework} sets up the light-cone Hamiltonian, the evolution operator $\Omega$, the scattering and the cross section, and recalls the LO result. Section~\ref{sec:lcwf} collects the LCWFs, and Sec.~\ref{sec:omega} the resulting coefficients of $\Omega$. Section~\ref{sec:xsec} gives the cross section at order $g^4$. The leading logarithms are extracted in Sec.~\ref{sec:LL}, the finite part in Sec.~\ref{sec:finite}, and the running coupling and DGLAP terms are discussed in Sec.~\ref{sec:rc}. We summarize in Sec.~\ref{sec:summary}. The appendices contain conventions, matrix elements and Fourier transforms (App.~\ref{app:conv}), the longitudinal integrals (App.~\ref{app:long}), the explicit form of JIMWLK$\otimes$LO and BFKL$\otimes$LO (App.~\ref{app:explicit}), the LL content of every term of the cross section (App.~\ref{app:LLrows}), and the numerical checks (App.~\ref{app:checks}).
